%% kapa: exam=1 round=22 year=2011 period=1 form=A subject=economics subject_ko=경제원론 seq=33 number=73
%% source: Q-Net 감정평가사 기출문제 (cst00309) artlSeq=5036790 file=extracted/1008/5036790/2011년 제22회 감정평가사 1차시험 기출문제/1교시.hwp
%% conversion: kordoc 4.24.51 -> tools/gen1.py (auto); review=proofread-sonnet; flags=none
\begin{question}{73}
다음은 국가간 자본이동이 완전한 소규모개방경제의 모형이다. 해외 이자율이 10으로 항상 일정할 때 중앙은행이 통화량을 20만큼 증가시킨다면, 통화량 증가 전과 후의 균형국민소득의 차이는? (단, Y는 국민소득, C는 소비, I는 투자, G는 정부지출, X는 수출, M은 수입, L\textsuperscript{D}는 실질화폐수요, L\textsuperscript{S}는 명목화폐공급, P는 물가수준, r은 국내 이자율, 국내 이자율 수준은 해외 이자율 수준과 같다.)
\begin{framedtext}
Y=C+I+G+X-M(생산물 시장의 균형) \par
$L^{ D } = \frac{ L^{ S } } { P }$(화폐시장의 균형) \par
C=50+0.8Y, I=100-4r, G=50 \par
X=30, M=30 \par
L\textsuperscript{D}=0.2Y-2r, L\textsuperscript{S}=100
\end{framedtext}
\begin{choices}
  \item 0
  \item 20
  \item 40
  \item 80
  \item 100
\end{choices}
\end{question}
\answer{1}
